\chapter{The 3D Viewport}
\label{ch:viewport}

The viewport is an OpenGL 3.3 core-profile canvas using instanced rendering
--- one draw call per mesh type --- so structures with tens of thousands of
atoms stay interactive.

\section{Interaction modes}
\label{sec:modes}

What a left-drag does depends on the active \emph{interaction mode}. The
six modes are the first group of buttons on the frame toolbar, each with a
single-letter shortcut:

\begin{center}
\small
\begin{tabular}{@{}llp{8.4cm}@{}}
\toprule
Key & Mode & Left mouse button \\
\midrule
\kbd{R} & Rotation   & Drag orbits the camera around the structure.
                       The default. \\
\kbd{T} & Translation & Drag pans the scene. \\
\kbd{S} & Selection  & Drag draws a rubber-band box and selects every atom
                       inside it. \\
\kbd{I} & Insertion  & Click empty space to add an atom; drag from one atom
                       to another to bond them. \\
\kbd{D} & Distance   & Click two atoms to measure their separation. \\
\kbd{A} & Angle      & Click three atoms (vertex second) to measure the
                       angle. \\
\bottomrule
\end{tabular}
\end{center}

Three gestures work in \emph{every} mode, so navigation never requires
switching back:

\begin{description}
  \item[Middle-drag, or \kbd{Shift}+left-drag] Pan.
  \item[Wheel] Zoom.
  \item[Double-click] Reframe the structure.
\end{description}

The mouse cursor changes with the mode --- an arrow for rotation, an open
hand for panning, a crosshair for selection, insertion and measurement.

\section{Selecting atoms}

A left \emph{click} (no drag) picks the atom under the cursor by ray-cast,
in any mode. \kbd{Ctrl}+click (\kbd{Cmd}+click on macOS) toggles an atom in
or out of the selection, so you can build up a set one atom at a time.
Clicking empty space clears the selection.

In \textbf{Selection} mode a rubber-band drag selects every atom whose
centre falls inside the box; holding \kbd{Ctrl} adds the boxed atoms to the
existing selection instead of replacing it. With the viewport focused,
\kbd{Delete} or \kbd{Backspace} removes the selected atoms and the bond
network is rebuilt immediately.

The status bar reports the selection size, and several tools act on it: the
bond-order buttons need exactly two atoms, the bond editor can pull a pair
straight from the selection, and \menu{Edit}\menusep\menu{Selection} can
change the element of, or translate, whatever is selected.

\section{Measuring distances and angles}
\label{sec:measure}

Measurement modes accumulate clicks on atoms and draw the result directly
on the canvas --- markers on the chosen atoms, dashed connecting lines, and
the value in a floating label. The same value is written to the status bar,
for example:

\begin{lstlisting}
Distance Si(0) - Si(4): 2.351 A
Angle O(2) - Si(0) - O(5): 109.47 deg
\end{lstlisting}

\begin{description}
  \item[Distance] Two atoms; result in \AA{} to three decimals.
  \item[Angle] Three atoms; the \emph{second} atom clicked is the vertex.
    Result in degrees to two decimals.
\end{description}

Clicking empty space clears the current measurement, and clicking again
after a completed one starts fresh. Dragging still orbits the camera, so
you can rotate the structure between clicks to reach an atom hidden behind
another. Measurements are cleared automatically when you switch modes or
load a different structure.

\section{Fixed-angle rotations}

The toolbar's rotation cluster turns the scene by exact increments about
the world axes --- useful for reproducible figure orientations.

\begin{description}
  \item[Angle step] An editable spin box, 1--180\textdegree, default
    15\textdegree.
  \item[\ui{X+} \ui{X$-$} \ui{Y+} \ui{Y$-$} \ui{Z+} \ui{Z$-$}]
    Counter-clockwise and clockwise rotation about each axis by exactly
    that step.
\end{description}

Each click animates smoothly over about 200\,ms, and rapid clicks compose
exactly: four presses of \ui{Z+} at 90\textdegree{} return the scene to
where it started. The axes triad follows the rotation, and picking,
selection and measurements all stay consistent with what you see.

\section{Framing, alignment and projection}

\begin{description}
  \item[\kbd{F}] Frame the structure: centre it and back the camera off far
    enough to see all of it. Also on the toolbar and in
    \menu{View}\menusep\menu{Alignment}.
  \item[Align with XY / XZ / YZ] Look straight down the remaining
    Cartesian axis. Toolbar buttons, or
    \menu{View}\menusep\menu{Alignment}. Alignment clears any accumulated
    fixed-angle rotation, so the result is a true axis-aligned view.
  \item[\kbd{O}] Toggle orthographic versus perspective projection. The
    orthographic frustum is matched to the perspective field of view at the
    camera's target distance, so the structure keeps its apparent size
    across the switch --- and zoom keeps working in both.
\end{description}

\begin{caltip}
Orthographic projection is what you want for figures where lattice
directions must read as parallel: perspective foreshortening makes a
supercell look subtly sheared.
\end{caltip}

\section{Inserting atoms and bonds}

In \textbf{Insertion} mode (\kbd{I}) the toolbar's element button --- which
shows the active symbol, \texttt{C} by default --- selects what gets
placed. Click it to choose any element from the graphical periodic table.

\begin{description}
  \item[Click on empty space] Inserts one atom of the active element there.
    The atom lands on the plane through the camera's target point facing
    the viewer, so what you click is where it appears.
  \item[Drag from one atom to another] Creates a bond between them, added
    as a manual bond override.
  \item[Click on an existing atom] Selects it, as in any other mode.
\end{description}

Both operations are single undo steps.

\begin{calnote}
Because inserted atoms land on the camera-target plane, the depth of a new
atom depends on the current view. Place atoms from an aligned, orthographic
view when the exact coordinate matters, or insert roughly and then set
exact coordinates with \menu{Edit}\menusep\menu{Add Atom}, which takes
numeric $x$, $y$, $z$.
\end{calnote}

\section{Visual effects}
\label{sec:visual-effects}

\menu{View}\menusep\menu{Visual Effects} toggles the \ui{Visual Effects}
dock --- the scene stays interactive while you tune it. It is a five-tab
panel running from what lights the scene, through the scene itself, to the
passes that post-process the finished image: \ui{Light} (the studio
lighting rig, Chapter~\ref{ch:appearance}), \ui{Shadow}, \ui{Fog},
\ui{Blur} and \ui{SSAO}. Every effect defaults to off. The ground plane
used to sit here as a \ui{Floor} tab; it now lives in the \ui{Spatial
References} dock instead (Chapter~\ref{ch:appearance}), since it draws a
whole extra object into the scene rather than an image filter.

\subsection{Shadows}

Real-time shadow mapping with percentage-closer filtering: atoms and bonds
cast shadows onto neighbouring geometry, including the floor when it is on.
Off by default; adds roughly one extra depth-only pass per frame.

\begin{description}
  \item[\ui{Intensity}] 0.00--1.00. How much direct light an occluded
    surface loses --- ambient light is never attenuated, so even at 1.00
    shadowed geometry stays readable rather than going black.
  \item[\ui{Softness / blur radius}] 0--6 texels. The PCF blur radius: 0
    gives hard, aliased edges; each step averages a wider neighbourhood, so
    quality and cost both rise together. 2--3 suits most structures.
\end{description}

The shadow projection follows the \textbf{primary light} --- the first
entry under the \ui{Light} tab. Re-aiming that light re-aims the shadows;
fill lights stay unshadowed, which is what keeps occluded regions legible.

\subsection{Distance fog}

Fades distant geometry into the background colour, which makes the depth
ordering of a thick slab or a large nanoparticle immediately readable.

\begin{description}
  \item[\ui{Mode}] \ui{Linear (start $\rightarrow$ end)} fades uniformly
    between two distances; \ui{Exponential (density)} follows
    $e^{-\rho d}$.
  \item[\ui{Start distance}, \ui{End distance}] The linear ramp, in \AA{}.
    Defaults 15 and 80.
  \item[\ui{Density}] The exponential coefficient, default 0.03.
\end{description}

The fog colour tracks the viewport background automatically, so geometry
always fades into whatever backdrop you have chosen.

\subsection{Depth of field}

A post-processing pass that blurs geometry away from the focal plane, in
proportion to its circle of confusion.

\begin{description}
  \item[\ui{Blur strength}] Maximum blur radius in pixels, 1--20.
  \item[\ui{Focus range}] The depth band that stays sharp, in \AA{}.
  \item[\ui{Focus offset}] Shifts the focal plane relative to the camera
    target, in \AA{}.
\end{description}

\begin{calnote}
Depth of field renders the scene into an offscreen buffer before
compositing, which means multisample anti-aliasing is unavailable while it
is enabled. For final figures, decide which of the two matters more ---
crisp edges or the depth cue.
\end{calnote}
